\section{Lucas's complete regional nonexistence argument}
\label{annex:sec:p18-lucas-no-solution}

\begin{annexitem}{annex:p18.lucas.region.partition}{thm:p18-lucas-region-partition}{All regions and the exhaustive residual cases}
\begin{proof}
The exact core formula is available from
\cref{thm:p18-lucas-core-foundation}. A point of \(E_i\) has
\[
                    x(O)=2+x_i+x_7+x_9<4.
\]
Thus \(E\cap C=\varnothing\); every defining \(F\) piece
excludes \(C\) explicitly. A point cannot be in two distinct
\(E_i\): its deficient first odd coordinate is unique, since
the other two are equal to 1.

Every point in \(F_{35},F_{51},F_{13}\) has \(x_7+x_9\geq1\).
This inequality also holds in \(F_7,F_9\), because one of
these coordinates is 1 and the other is nonnegative, and in
\(F_{79}\), where their sum is 2. Hence none of these pieces
meets \(E\). In \(F_{135}\) the first three odd coordinates
are all 1, whereas each \(E_i\) requires one of them to be
less than 1. This proves \(E\cap F=\varnothing\). The set
\(R\) is defined to remove exactly \(C,E,F\) from \(B\);
all four are subsets of \(B\). Adjoining \(A\setminus B\)
therefore gives the claimed disjoint partition.

Pair sums in \(B\) force every coordinate to lie in \([0,1]\),
and determine the even coordinates from the odd ones.
The formula for \(e_i\) consequently describes every point
of \(E_i\) and only such points. The values \(t=u=v=0\)
show nonemptiness, including the permitted lower boundaries.

Let \(x\in R\). Since \(x\notin C\), its odd sum is less
than 4. If all three first odd coordinates were 1, then
\(x\in F_{135}\), a contradiction. If exactly two, say
\(x_j=x_k=1\), were 1, the third is less than 1. The two
possibilities \(x_7+x_9<1\) and \(x_7+x_9\geq1\) would put
\(x\) in \(E_i\) and \(F_{jk}\), respectively. Thus at most
one first odd coordinate equals 1.

Both \(x_7=x_9=1\) would put \(x\) in \(F_{79}\), so at
least one is less than 1. If \(x_p=1\), the other coordinate
satisfies \(x_q<1\). Exclusion from \(F_p\), together with
exclusion from \(C\), then forces at least one of the three
displayed pair-plus-\(q\) sums to be strictly less than 2.
Equality in all three would satisfy the definition of
\(F_p\), so no equality case was discarded.
\end{proof}
\end{annexitem}

\begin{annexitem}{annex:p18.lucas.core.outside}{thm:p18-lucas-core-outside}{A core witness for every deficient pair}
\begin{proof}
The five pair sums of a point in \(A\) add to 5. If all
were at least 1, they would all be 1 and the point would
belong to \(B\). Thus a point outside \(B\) has a pair
\(\{h-1,h\}\) of sum less than 1.
The stated \(\epsilon\) is positive, and the two coefficients
of \(c^0,c^h\) are nonnegative with sum
\[
                     x_{h-1}+x_h+2\epsilon=1.
\]
Both vertices are in \(C\), and the core is their containing
convex simplex, so \(y\in C\). At the selected pair,
\[
                  y_{h-1}=x_{h-1}+\epsilon,\qquad
                  y_h=x_h+\epsilon.
\]
Both improvements are strict, while their sum is exactly
1, the pair's coalition value. Hence
\(y\succ_{\{h-1,h\}}x\). This argument permits either original
coordinate to be zero; it uses no strict positivity of an
imputation's coordinates.
\end{proof}
\end{annexitem}

\begin{annexitem}{annex:p18.lucas.core.remainder}{thm:p18-lucas-core-remainder}{Every residual coordinate case and its feasible coalition}
\begin{proof}
First record the common construction. Suppose \(T\subseteq O\)
has \(m\) members, all their coordinates are less than 1,
and \(s=x(T)<m-1\). Then
\[
 0<m-1-s<m-s,\qquad
 0<\lambda=\frac{m-1-s}{m-s}<1.
\]
For \(h\in T\), the proposed \(y_h\) is strictly greater
than \(x_h\), less than 1, and nonnegative. Its sum is
\[
                \sum_{h\in T}y_h=s+\lambda(m-s)=m-1.
\]
Set the other odd coordinates to 1 and complete the even
coordinates by pairing. This gives a point in \(B\) with
odd sum \((m-1)+(5-m)=4\), hence in \(C\).
More explicitly, its simplex coordinates are
\[
                      y=\sum_{h\in T}(1-y_h)c^{h+1}.
\]
The coefficients are positive and sum to 1. This identifies
the exact core face providing the dominator.

Now take \(x\in R\). If \(x_7<1\) and \(x_9<1\), there
are two exhaustive possibilities. When none of
\(x_1,x_3,x_5\) is 1, use \(T=O\), \(m=5\);
its sum is less than 4 by the region lemma. The witness
satisfies \(y(O)=4=\ell(O)\) and strictly improves all five
odd players, so it dominates via \(O\).

When one first odd coordinate equals 1, name it \(x_i\).
The other two \(x_j,x_k\) are less than 1. With
\(T=\{j,k,7,9\}\), all four selected coordinates are less
than 1, and
\[
                        x(T)=x(O)-1<3.
\]
Use \(m=4\), leave \(y_i=1\), and obtain
\(y(T)=3=\ell(O\setminus\{i\})\). This is domination by
that four-member odd coalition.

In the remaining possibility exactly one of \(x_7,x_9\)
equals 1. Name it \(x_p\), so \(x_q<1\). Choose a cyclic
triple \((i,j,k)\) with \(x_j<1\) and \(x_k<1\).
Such a choice is possible because at most one first odd
coordinate equals 1; if there is one, assign its index to
\(i\), and if there is none, any cyclic triple works.
At least one of the three pair-plus-\(q\) sums is less
than 2.

If \(x(jkq)<2\), take \(T=\{j,k,q\}\).
All its coordinates are less than 1, so the construction
with \(m=3\) applies. The two omitted odd coordinates
\(y_i,y_p\) are set to 1, and
\[
                y(T)=2=\ell(O\setminus\{i,p\}).
\]
Suppose instead that a deficient triple is
\(\{i,h,q\}\), with \(h\in\{j,k\}\).
If \(x_i<1\), take this triple itself for \(T\), set the
remaining first odd coordinate and \(y_p\) to 1, and use
the corresponding value-2 coalition.
If \(x_i=1\), let \(n\) be the other member of \(\{j,k\}\).
Then \(x_n<1=x_i\), and
\[
                x(jkq)=x_h+x_n+x_q
                         <x_h+x_i+x_q=x(ihq)<2.
\]
Thus the earlier selection \(T=\{j,k,q\}\) applies.
This handles both deficient triples \(\{i,j,q\}\) and
\(\{i,k,q\}\), both orientations of \((p,q)\), and all
three cyclic triples.

These cases exhaust \(R\). They always select a coalition
\(T\) of its exact source value \(|T|-1\), and all its members
strictly improve. The preceding outside-cube lemma now
gives the full displayed containment in \(\operatorname{Dom}C\).
\end{proof}
\end{annexitem}

\begin{annexitem}{annex:p18.lucas.incoming.regions}{thm:p18-lucas-incoming-regions}{Exhausting every coalition and every exceptional face}
\begin{proof}
Several exclusions apply before any regional calculation.
No imputation dominates another through \(N\): strict
improvement of all ten coordinates contradicts equality
of their total sums. A zero-value nonempty coalition cannot
support domination, because all target coordinates are
nonnegative, strict improvement makes the attacker's
coalition sum positive, and feasibility would require it
to be at most zero. No point of \(B\) can be dominated
through a paired coalition by any point of \(A\):
strictly improving both coordinates would raise their sum
above its value 1.

Take a target \(x\in E_i\) and an attacker \(y\in B\).
Both \(x_j,x_k\) are 1, while every coordinate of \(y\)
is at most 1. No attacking coalition can contain \(j\)
or \(k\). Apart from the pairs and \(N\) already excluded,
the positive-value coalitions are exactly \(O\), its three
listed four-member subsets, its six listed three-member
subsets, and \(M_1,M_3,M_5\).
The first three types each contain at least one of \(j,k\):
they omit at most one member of \(\{1,3,5\}\).
Among the mixed coalitions only \(M_i\) avoids both \(j\)
and \(k\). This proves the claimed unique possible coalition.
Moreover \(x_r=x_{j+1}=0\), so any such attacker must
have \(y_r>0\).

We next prove that a point in \(C\cup F\) feasible through
\(M_i\) always has \(y_r=0\).
If \(y\in C\), use \(y_r=1-y_j\) and \(y(O)\geq4\):
\[
             y(M_i)=y(O)+1-2y_j-y_k
                                      \geq5-2y_j-y_k.
\]
If this is at most 2, then \(2y_j+y_k\geq3\).
Since \(y_j,y_k\leq1\), both must be 1; in particular
\(y_r=0\).

For \(y\in F\), examine every component. If \(y\in F_{79}\),
then \(y_7+y_9=2\), so feasibility gives
\(y_i+y_r\leq0\), forcing \(y_r=0\).
If \(y\in F_{135}\), all three first odd coordinates are
1 and \(y_r=1-y_j=0\). If \(y\in F_{jk}\) or \(F_{ij}\),
again \(y_j=1\), so \(y_r=0\).
For the remaining first-odd pair face \(F_{ik}\), one has
\(y_i=1\) and \(y_7+y_9\geq1\); hence
\[
                     y(M_i)=1+y_r+y_7+y_9\geq2+y_r.
\]
Feasibility forces \(y_r=0\).
Finally let \(y\in F_p\), for either orientation of
\((p,q)\). Its inequality for the pair \(\{i,k\}\) gives
\(y_i+y_k+y_q\geq2\). Since \(y_k\leq1\),
\(y_i+y_q\geq1\). Together with \(y_p=1\), this yields
\[
                     y(M_i)=y_i+y_q+1+y_r\geq2+y_r,
\]
again forcing \(y_r=0\).
All components have been covered, including their overlaps.
The result contradicts the necessary strict improvement
at \(r\) for any target in \(E_i\), proving
\(E\cap\operatorname{Dom}(C\cup F)=\varnothing\).

For completeness, prove the source's additional assertion
about targets in \(F\). Let \(x\in F\) and
\(y\in C\cup E\cup F\subseteq B\). The exclusions for
\(N\), zero-value coalitions and pairs still apply.
If \(x\in F_{79}\), every remaining positive coalition
contains at least one of players 7 and 9, whose target
coordinate is 1; domination is impossible.
If \(x\in F_{135}\), every remaining positive coalition
contains a first odd player, again with coordinate 1.

If \(x\in F_p\), all remaining positive coalitions containing
\(p\) are blocked by \(x_p=1\).
The only remaining positive coalitions without \(p\) are
the three triples \(O\setminus\{h,p\}\), \(h\in H\).
Their target sums are at least 2 by the definition of
\(F_p\). Strict improvement of their members would give
an attacker sum greater than their value 2. These also
cannot dominate.

It remains to consider \(x\in F_{jk}\), for each of the
three first-odd pair faces. Its coordinates \(x_j=x_k=1\)
exclude every positive odd coalition and all mixed
coalitions except \(M_i\), just as above. An attacker
through \(M_i\) requires \(y_r>0\).
If \(y\in C\cup F\), the already proved feasibility
calculation rules this out. If \(y\in E_i\cup E_k\),
then \(y_j=1\) and \(y_r=0\). If \(y\in E_j\), then
\(y_7+y_9<1\), whereas strict improvement at both 7 and
9 would require
\[
                          y_7+y_9>x_7+x_9\geq1.
\]
This is impossible as well. We have considered every
target piece and every possible attacker, proving the
second immunity formula. All arguments use weak
feasibility and retain equality on the stated faces.
\end{proof}
\end{annexitem}

\begin{annexitem}{annex:p18.lucas.cyclic.dominance}{thm:p18-lucas-cyclic-dominance}{The exact attack relation on the three wedges}
\begin{proof}
For a target \(x\in E_i\), the only possible coalition
for an attacker in \(E\subseteq B\) is \(M_i\).
If \(y\in E_i\cup E_k\), then \(y_j=1\), so
\(y_{j+1}=0=x_{j+1}\); strict improvement fails.
Thus the attacker must be in \(E_j\).

Write \(y=e_j(t,u,v)\). Its coordinates on \(M_i\) are
\[
                 y_i=1,\quad y_{j+1}=1-t,\quad
                 y_7=u,\quad y_9=v.
\]
For the target they are \(x_i<1\), \(x_{j+1}=0\),
\(x_7,x_9\). Because \(t<1\), the first two improvements
are automatic and strict. The other two are precisely
\(u>x_7\) and \(v>x_9\). Feasibility is exactly
\[
              1+(1-t)+u+v\leq2
                       \quad\Longleftrightarrow\quad t\geq u+v.
\]
These necessary conditions are sufficient, since they
verify every member's improvement and the whole coalition
inequality. Cycling the indices gives \(E_3\) attacking
\(E_1\), \(E_5\) attacking \(E_3\), and \(E_1\) attacking
\(E_5\), with no other directions.

If \(t<u+v\), the one possible outgoing coalition is
infeasible. At equality it is feasible; the coordinate
inequalities alone determine its targets. If \(u=0\),
there is no target with \(x_7<u\), since \(x_7\geq0\);
the same argument applies to \(v=0\). Thus feasibility
alone need not provide an outgoing target. Conversely,
when \(u,v>0\), points of \(E_i\) with, for example,
last odd coordinates \(u/2,v/2\) provide targets whenever
the feasibility inequality holds. Their sums are less
than 1 because \(u+v<1\).
\end{proof}
\end{annexitem}

\begin{annexitem}{annex:p18.lucas.upper.regions}{thm:p18-lucas-upper-regions}{Explicit witnesses for every geometric assertion}
\begin{proof}
Let \(s=z_7+z_9<1\). With
\[
 \delta=\frac{1-s}{4},\qquad
 u=z_7+\delta,\quad v=z_9+\delta,\quad t=u+v,
\]
one has \(u>z_7\), \(v>z_9\), and
\[
                         t=s+2\delta=\frac{1+s}{2}<1.
\]
Consequently \(e_h(t,u,v)\in G_h(z)\) for every \(h\).
This also works when \(z_7=0\), \(z_9=0\), or the
active first odd coordinate of \(z\) is zero.

We prove the general identity
\[
                     \operatorname{Pred}_E(G_h(z))
                                               =G_{\rho(h)}(z).
\]
If \(w\) dominates \(x\in G_h(z)\), the exact cyclic
relation forces \(w\in E_{\rho(h)}\), with its active
coordinate at least \(w_7+w_9\). Also
\[
                        w_7>x_7>z_7,\qquad
                        w_9>x_9>z_9.
\]
Hence \(w\in G_{\rho(h)}(z)\).

Conversely take \(w\in G_{\rho(h)}(z)\), and put
\[
             a=\frac{z_7+w_7}{2},\qquad
             b=\frac{z_9+w_9}{2},\qquad x=e_h(a+b,a,b).
\]
The inequalities \(z_7<a<w_7\) and \(z_9<b<w_9\) are strict.
Moreover \(a+b<w_7+w_9<1\), so \(x\in G_h(z)\).
The point \(w\) satisfies its outgoing feasibility
inequality by membership in \(G_{\rho(h)}(z)\);
the strict inequalities at 7 and 9 now give \(w\succ x\).
This proves the reverse inclusion. Substitution of
\(h=i,j,k\) supplies all three predecessor identities,
including all orientations of the cycle.

Next take \(x=e_i(t,u,v)\in G_i(z)\). Since \(u+v<1\), set
\[
 \epsilon=\frac{1-u-v}{4},\qquad
 a=u+\epsilon,\quad b=v+\epsilon,\qquad
                         y=e_j(a+b,a,b).
\]
Then \(a+b=(1+u+v)/2<1\), both coordinates strictly
exceed those of \(x\), and hence those of \(z\).
Thus \(y\in G_j(z)\), with its outgoing coalition
feasible at equality, and \(y\succ x\).
This proves \(G_i(z)\subseteq\operatorname{Dom}_E(G_j(z))\).

Finally, for \(z\in E_j\), the exact cyclic domination
formula says that its attackers are the points in
\(E_k\) with both last odd coordinates strictly above
those of \(z\) and with active coordinate at least their
sum. This is exactly \(G_k(z)\).
All witnesses used actual points with active coordinate
less than 1 and coordinate sum less than 1. Thus none
rests on an unattained upper boundary or on a limit
assumption about an arbitrary candidate solution.
\end{proof}
\end{annexitem}

\begin{annexitem}{annex:p18.lucas.local.nonexistence}{thm:p18-lucas-local-nonexistence}{The complete internal--external stability contradiction}
\begin{proof}
Suppose \(K'\subseteq E\) were a relative solution.
The set \(E\) is nonempty, so \(K'\) cannot be empty:
an empty set dominates no outsider. Choose \(z\in K'\),
and fix the cyclic triple with \(z\in E_j\).
Internal stability, together with the exact predecessor
set of \(z\), gives
\[
                              K'\cap G_k(z)=\varnothing. \tag{L1}
\]
The set \(G_k(z)\) is nonempty. Choose any \(x\) in it.
By (L1), \(x\notin K'\), so external stability supplies
\(w\in K'\) with \(w\succ x\).
The predecessor identity
\(\operatorname{Pred}_E(G_k(z))=G_i(z)\) forces \(w\in G_i(z)\).
Therefore
\[
                              K'\cap G_i(z)\ne\varnothing. \tag{L2}
\]

Now take any \(y\in G_j(z)\). If \(y\notin K'\), external
stability would supply an attacker in
\[
                 K'\cap\operatorname{Pred}_E(G_j(z))
                                      =K'\cap G_k(z),
\]
contrary to (L1). Hence every such \(y\) belongs to \(K'\):
\[
                                         G_j(z)\subseteq K'. \tag{L3}
\]
Every point of \(G_i(z)\) is dominated by a point of \(G_j(z)\),
by the explicit witness lemma. In view of (L3), internal
stability now forces
\[
                              K'\cap G_i(z)=\varnothing,
\]
contradicting (L2). This establishes nonexistence for
arbitrary subsets \(K'\); no cardinality, topological,
or measure assumption entered the proof.
\end{proof}
\end{annexitem}

\begin{annexitem}{annex:p18.lucas.no.solution}{thm:p18-lucas-no-solution}{Lifting the regional contradiction to the whole game}
\begin{proof}
Suppose \(K\subseteq A\) were an NM solution.
By \cref{thm:p18-lucas-core-foundation}, the core is precisely
the undominated set. It must be
contained in every solution: an omitted undominated point
could not satisfy external stability. Thus \(C\subseteq K\).
Internal stability then excludes every point of
\(\operatorname{Dom}C\) from \(K\). The two core-domination
lemmas and the disjoint partition imply
\[
                                         K\subseteq C\cup E\cup F.
\]
Put \(K'=K\cap E\). It inherits internal stability.
For any \(x\in E\setminus K'\), one has \(x\notin K\), so
external stability of \(K\) provides \(y\in K\) with
\(y\succ x\). The inclusion just proved restricts \(y\)
to \(C\cup E\cup F\), and the incoming-region lemma excludes
\(y\in C\cup F\). Thus \(y\in K'\).
Consequently \(K'\) dominates all of \(E\setminus K'\);
it is a relative solution, contradicting the local
nonexistence theorem.

The additional \(F\)-immunity formula would also force
\(F\subseteq K\), since an omitted point of \(F\) has no
attacker in the only region where \(K\) could lie.
The contradiction above already follows without this
additional consequence.

The argumentation framework uses all of \(A\) and exactly
this domination relation as its attack relation.
Being conflict-free is exactly internal stability, and
attacking every outside argument is exactly external
stability. Thus its stable extensions would be precisely
the NM solutions just excluded. This direct equivalence
uses the definitions and needs no superadditivity assumption
on \(\ell\). Alternatively, the already proved foundation
identifies this graph with that of the superadditive cover,
where \cref{thm:p18-nm-argumentation} applies directly.
The previously established core simplex remains nonempty;
for example \(c^0\in C\). None of its construction has
been altered by the regional argument.
\end{proof}
\end{annexitem}

\begin{annexitem}{annex:p18.lucas.positive.cover}{thm:p18-lucas-positive-cover}{Every positive coalition under the partition cover}
\begin{proof}
We verify the source's assertion for every positive
coalition, including those not required merely to
preserve the grand-coalition value.

For \(N\), every listed value satisfies
\(\ell(S)\leq|S\cap O|\), and unlisted values are zero.
Summing over a partition of \(N\) therefore gives a
total at most 5. The one-block partition attains 5.
Thus \(\widehat\ell(N)=5=\ell(N)\).
For a paired coalition \(Q_a\), the same bound gives
every partition a total at most 1, attained by the
one-block partition.

Consider now a coalition consisting exclusively of odd
players. Its positive subcoalitions can only be \(O\)
(value 4), one of the listed odd four-member sets
(value 3), or one of the listed odd three-member sets
(value 2). A pair \(Q_a\) or a mixed \(M_i\) cannot be
contained in such a coalition, because each has an even
player. Each positive odd subcoalition has at least
three members, so two of them cannot be disjoint inside
a set of at most five odd players. Thus a partition
of any positive odd coalition has at most one
positive-value block.
For \(O\), that block has value at most 4; for a listed
four-member coalition it has value at most 3; for a
listed three-member coalition it has value at most 2.
The respective one-block partitions attain these bounds.
This proves equality of the cover at all ten positive
odd coalitions.

Finally take \(M_i=\{i,j+1,7,9\}\).
It has only one even player, \(j+1\), and its paired
odd player \(j\) is absent, so it contains no \(Q_a\).
Its odd part \(\{i,7,9\}\) contains only one member of
\(H\), whereas every positive odd triple in the game
contains two members of \(H\). It therefore contains
no positive odd subcoalition. A mixed coalition
contained in \(M_i\) must have the same unique even
player \(j+1\); among the three mixed coalitions this
uniquely selects \(M_i\) itself. Nor can \(N\) be a
subcoalition of \(M_i\). Consequently every proper
block in a partition of \(M_i\) has value zero, while
the one-block partition has value 2. This proves
\(\widehat\ell(M_i)=2\) for all three cyclic indices.

The list now covers every coalition with positive
\(\ell\)-value. Singleton values are zero and any singleton
has just its one-block partition, so those values are
unchanged. By \cref{thm:p18-superadditive-cover}, the cover
is superadditive and equality at \(N\) makes the imputation
sets and entire domination graphs identical.
For clarity, in the nontrivial direction an attack
feasible under \(\widehat\ell\) has a maximizing partition
of its coalition. At least one block is feasible under
\(\ell\); otherwise summing strict violations over all
blocks would contradict cover feasibility. Every member
of that block still strictly improves. Thus it was
already an attack in the original graph. The reverse
direction follows from \(\widehat\ell\geq\ell\).
The core and solution families are therefore identical,
and the proved nonexistence transfers to
\((N,\widehat\ell)\).
\end{proof}
\end{annexitem}
